\subsection*{Neutrality and Interpretation}
\textit{Terms specific to the neutrality requirements defined in this paper.}

%-----------------------------------------------------
\subsubsection*{Neutrality Requirements}
%-----------------------------------------------------

\begin{description}[style=nextline,leftmargin=1.5cm]

\item[Neutrality]
The property of a substrate ontology that satisfies
both interpretive non-commitment and extension stability.
Neutrality, in this paper, is not moral neutrality,
political impartiality, epistemic skepticism,
or absence of content.
It is a role-relative structural property:
the substrate must provide stable shared reference
without asserting framework-variant conclusions
as substrate-layer facts.

\item[Interpretive Non-Commitment]
The requirement that a substrate not assert
framework-variant propositions as substrate-layer commitments.
In particular, a neutral substrate must not assert
causal or normative conclusions as substrate-layer facts
when their truth depends on an admissible interpretive framework.
Interpretive non-commitment does not prevent the substrate
from representing that such claims were made;
it prevents the substrate from asserting those claims as true.
Together with \emph{Extension Stability}, this requirement
defines \emph{Neutrality}.

\item[Extension Stability]
The requirement that a substrate remain logically consistent
when extended separately by each admissible interpretive framework.
Extension stability is pairwise:
for each admissible framework $\mathcal{F}$,
the combination $\mathcal{S} \cup \mathcal{F}$
must remain consistent.
It does not require all admissible frameworks
to be mutually consistent with one another.
Together with \emph{Interpretive Non-Commitment},
this requirement defines \emph{Neutrality}.

\item[Neutral Substrate]
A substrate ontology that satisfies neutrality.
A neutral substrate provides stable shared reference
across admissible interpretive frameworks
without asserting framework-variant conclusions
as substrate-layer facts.
It may represent causal or normative claims
as reified, attributed, provenance-bearing content,
but it does not assert those claims as true
at the substrate layer.

\item[Substrate-Layer Fact]
A proposition asserted by the substrate ontology
as holding independently of any particular interpretive framework.
Substrate-layer facts are part of the shared base
over which admissible frameworks may extend,
reason, disagree, or derive conclusions.
In a neutral substrate,
substrate-layer facts must be restricted to
framework-invariant referential structure
and other framework-invariant structural commitments.

\item[Framework-Level Conclusion]
A proposition derived within an interpretive framework
from substrate-layer facts together with
the framework's own rules, assumptions, methods, or norms.
Framework-level conclusions may vary across admissible frameworks.
They are not asserted by the substrate merely because
they can be derived by some framework.

\end{description}

%-----------------------------------------------------
\subsubsection*{Interpretive Frameworks}
%-----------------------------------------------------

\begin{description}[style=nextline,leftmargin=1.5cm]

\item[Interpretive Framework]
A coherent set of principles, rules, methods, assumptions,
or norms used to evaluate, explain, classify,
or assign meaning to shared referents.
Examples include legal jurisdictions, causal modeling traditions,
ethical theories, regulatory schemes, and analytic methodologies.
Interpretive frameworks may disagree with one another
while each remains internally consistent.

\item[Admissible Framework]
An interpretive framework that is internally consistent
and recognized as a legitimate stance
within the domain of application.
Admissibility is treated in this paper
as a domain-governance parameter,
not as something determined by the substrate itself.
Admissible frameworks need not be mutually compatible.

\item[Pluralism]
The coexistence of multiple,
potentially incompatible interpretive frameworks
within a single domain or system.
A neutral substrate supports pluralism
by providing shared referents and framework-invariant structure
without committing the substrate to the conclusions
of any single framework.

\item[Accountability System]
A data system designed to support explanation,
evaluation, contestation, or assignment of responsibility
across multiple stakeholders or interpretive perspectives.
Accountability systems are a motivating application
for neutral substrates because they often must accommodate
persistent disagreement about causes, obligations,
interpretations, and responsibility.
This paper uses accountability as an application context;
it does not define a specific accountable-record architecture.

\end{description}

%-----------------------------------------------------
\subsubsection*{Commitment and Derivation}
%-----------------------------------------------------

\begin{description}[style=nextline,leftmargin=1.5cm]

\item[Commitment]
A proposition asserted by an ontology as holding
at the substrate layer,
independently of any particular interpretive framework.
Commitments are the content of ontological assertions:
what the ontology treats as true,
not merely what has been claimed by external agents,
documents, institutions, or frameworks.
In a neutral substrate,
framework-variant propositions may not be asserted
as substrate-layer commitments.
Contrast with \emph{Conclusion}.

\item[Conclusion]
A proposition derived by an interpretive framework
from substrate-layer commitments together with
the framework's own rules, assumptions, methods, or norms.
Conclusions belong to interpretive layers,
not to the substrate.
Contrast with \emph{Commitment}.

\item[Externalization]
The architectural practice of locating contested interpretation
in interpretive layers above the substrate,
rather than embedding contested conclusions
as substrate-layer commitments.
Externalization preserves neutrality by ensuring that
framework-dependent conclusions remain associated
with the frameworks that derive or assert them.
Reification is one mechanism for externalization.

\item[Substrate Revision]
The modification, retraction, or replacement
of substrate-layer commitments in response to
the introduction of a new interpretive framework.
Substrate revision violates extension stability
when it is required in order to accommodate
an admissible framework.
Under the neutrality definition,
interpretive conflict must not be resolved
by revising the substrate's framework-invariant commitments.

\end{description}

%-----------------------------------------------------
\subsubsection*{Compatibility and Scope}
%-----------------------------------------------------

\begin{description}[style=nextline,leftmargin=1.5cm]

\item[Consistency ($\not\vdash \bot$)]
The property of a set of statements such that
no contradiction can be derived from its assertions.
The notation $\mathcal{S} \not\vdash \bot$
indicates that a substrate ontology $\mathcal{S}$
is consistent.
The notation $\mathcal{S} \vdash \bot$
indicates that $\mathcal{S}$ is inconsistent.

\item[Compatibility]
The property of a substrate ontology
and an interpretive framework such that
their combination does not entail contradiction.
Formally, $\mathcal{S} \cup \mathcal{F} \not\vdash \bot$.
In this paper, compatibility is assessed pairwise
between the substrate and each admissible framework.

\item[Pairwise Compatibility]
The requirement that each admissible framework
be separately compatible with the substrate.
Pairwise compatibility does not require admissible frameworks
to be compatible with one another.
This distinction allows the substrate to support
mutually incompatible interpretive frameworks
without collapsing their disagreement.

\item[Boundary Condition]
A precondition that must hold for a framework or result to apply.
In this paper, the central boundary condition is that
the substrate's referential regime is fixed at the substrate layer
and remains invariant
across the admissible frameworks under consideration.
Domains where individuation, co-reference, or persistence
are themselves contested do not admit neutral substrates
in the sense defined here until that contestation is resolved
or scoped out.
Boundary conditions delimit scope;
they are not defects but explicit constraints on applicability.

\item[Scope Condition]
A condition limiting the domain over which
the neutrality theorem applies.
The central scope condition in this paper is
referential invariance:
the relevant referential regime must be fixed
at the substrate layer
and remain invariant
across the admissible frameworks under consideration.
If the referential regime itself is contested,
the theorem does not claim that a neutral substrate
is available until the contestation is resolved
or excluded from the domain of application.

\end{description}

%-----------------------------------------------------
\subsubsection*{Invariant and Variant Content}
%-----------------------------------------------------

\begin{description}[style=nextline,leftmargin=1.5cm]

\item[Framework-Invariant Referential Structure]
The referential structure whose interpretation
remains stable
across the admissible frameworks under consideration.
In this paper, framework-invariant referential structure
includes the substrate's fixed referential regime:
the individuation, co-reference, and persistence conditions
by which entities, occurrences, and institutional artifacts
are identified and tracked.
Such structure may be asserted at the substrate layer
because its interpretation does not depend on
a particular admissible framework.
Contrast with \emph{Framework-Variant Proposition}.

\item[Framework-Variant Proposition]
A proposition whose truth value may differ across
admissible interpretive frameworks.
A framework-variant proposition may be derived,
asserted, or evaluated within an interpretive framework,
but it may not be asserted as a substrate-layer commitment
under the neutrality definition.
Formally, framework variance is defined in
Definition~\ref{def:framework-variant}.
Contrast with \emph{Framework-Invariant Referential Structure}.

\item[Shared Referent]
An entity, occurrence, or institutional artifact
to which multiple interpretive frameworks refer
through the substrate's referential regime.
Shared referents provide common ground:
they allow incompatible interpretive conclusions
to be recognized as conclusions about the same thing,
rather than as merely equivocal uses of different references.

\end{description}

%-----------------------------------------------------
\subsubsection*{Disagreement}
%-----------------------------------------------------

\begin{description}[style=nextline,leftmargin=1.5cm]

\item[Persistent Interpretive Disagreement]
The condition in which multiple admissible interpretive frameworks
reach incompatible conclusions about shared referents,
and the disagreement is structural rather than merely
a temporary result of missing data, translation failure,
or lack of coordination.
Persistent interpretive disagreement is the motivating condition
for neutral substrate design.

\item[Semantic Heterogeneity]
The condition in which multiple ontologies,
schemas, or vocabularies exist for the same domain,
and disagreement is treated primarily as a translation,
alignment, mapping, or mediation problem.
Semantic heterogeneity often assumes that disagreement
can be resolved through better alignment.
Persistent interpretive disagreement, by contrast,
may remain even when reference is shared
and terminology is aligned.
Contrast with \emph{Persistent Interpretive Disagreement}.

\end{description}
